% --- Corps de chapitre à inclure : % --- Corps de chapitre à inclure : % --- Corps de chapitre à inclure : % --- Corps de chapitre à inclure : \input{docs/module_conception}
% Si votre classe est article : remplacer \chapter par \section{Conception}
% Nécessite : \usepackage{listings} (ou minted)

\chapter{Conception}

Dans ce chapitre, nous exposons la conception de la partie applicative et péristaltique du système de détection du churn : la manière dont nous structurons les données, la génération contrôlée de \emph{mock} pour le développement et les tests, le modèle métier \texttt{ClientChurn}, et l'articulation avec le reste de la chaîne (datasets réels, API de prédiction, interface). Nous adoptons un point de vue volontairement ancré dans le code : le professeur a insisté sur la présence d'extraits représentatifs ; nous en présentons donc plusieurs, tirés du dépôt tel qu'il existe aujourd'hui.

\section{Périmètre et objectifs de conception}

Nous concevons le système pour qu'un utilisateur rattaché à une \textbf{agence} ne voie et ne modifie que les clients et analyses de cette agence, conformément à la hiérarchie métier. Nous séparons l'interface web (Django), le calcul de scores côté service dédié (FastAPI) lorsqu'il est disponible, et la persistance relationnelle (PostgreSQL via Django). Nous prévoyons en outre un repli sur un scoring local lorsque l'API est injoignable, afin de ne pas bloquer une démonstration ou un environnement de test.

\section{Données mock : rôle, garde-fous et extrait de code}

Les données \emph{mock} nous permettent de peupler rapidement la base avec des profils clients cohérents (télécom, réclamations, usage) lorsque aucun fichier CSV réel n'est chargé, ou pour valider le pipeline bout en bout. Nous avons toutefois restreint explicitement ce mode : si l'application tourne avec \texttt{DEBUG=False} (production), nous interdisons l'appel à la génération mock pour éviter d'injecter des données fictives en situation réelle. Le principe est implémenté dans \path{core/mock_data.py} : une vérification au moment de l'exécution, et une levée d'exception si le mode n'est pas autorisé.

\begin{lstlisting}[language=Python, caption={Garde-fou avant génération mock (\path{core/mock_data.py}, extrait).}]
def _ensure_mock_mode():
    """Leve une exception si le mode Mock est utilise en production."""
    if not _check_mock_mode():
        raise RuntimeError(
            "ERREUR CRITIQUE: Le mode Mock (generer_mock_data) est desactive "
            "en production. DEBUG=False detecte."
        )
\end{lstlisting}

L'entrée principale \texttt{generer\_mock\_data(agence\_id=..., user\_id=..., nb\_clients=50)} résout l'agence et l'utilisateur, crée un \texttt{Dataset} avec \texttt{methode="mock"}, puis génère des enregistrements \texttt{ClientChurn} avec des tirages aléatoires contrôlés (plans tarifaires, CDR agrégés, règles d'audit RGS-90 comme la division protégée pour la durée moyenne d'appel). Nous illustrons ci-dessous la signature et le début de la fonction : c'est le cœur du branchement « données simulées » attendu en conception.

\begin{lstlisting}[language=Python, caption={Entrée de la génération mock (extrait).}]
def generer_mock_data(agence_id=None, user_id=None, nb_clients=50):
    """
    Genere des donnees mock pour les tests.
    SECURITE: bloquee en production (DEBUG=False).
    """
    _ensure_mock_mode()
    from learning.models import ClientChurn, Dataset
    from core.models import Agence
    # ... resolution agence / utilisateur, Dataset(methode="mock"), boucle clients
\end{lstlisting}

Par ailleurs, nous structurons des règles métier de recommandation sous forme de liste de dictionnaires (\texttt{REGLES\_RECOMMANDATIONS}) : chaque entrée combine une condition sur un objet \texttt{ClientChurn} (par exemple nombre de réclamations) et un texte de mission généré dynamiquement. Cette conception nous permet de documenter clairement le lien entre attributs du client mock ou réel et actions proposées dans les notifications, sans durcir ces textes dans la vue.

\section{Modèle \texttt{ClientChurn} et lien \texttt{Dataset} / agence}

Nous modélisons le client comme une entité unique \texttt{ClientChurn} portant à la fois les informations de profil (identité, contrat, usage agrégé) et les sorties du modèle prédictif (\texttt{score\_churn}, \texttt{churn\_predit}, \texttt{niveau\_risque}). Conformément aux migrations du projet, chaque client est rattaché à une \textbf{agence} par clé étrangère, et optionnellement à un \texttt{Dataset} lorsque les lignes proviennent d'un import (CSV) ou d'un jeu mock. Nous complétons l'agence à la sauvegarde si elle manque mais qu'un dataset est présent, pour respecter la contrainte de base et les filtres par agence dans les vues.

\begin{lstlisting}[language=Python, caption={Structure essentielle de \texttt{ClientChurn} (\path{learning/models.py}, extrait).}]
class ClientChurn(models.Model):
    agence = models.ForeignKey("core.Agence", on_delete=models.CASCADE)
    dataset = models.ForeignKey(
        "Dataset", on_delete=models.CASCADE,
        related_name="clients", null=True, blank=True,
    )
    client_id = models.CharField(max_length=100)
    # ... profil, usage (nb_appels, data_totale_mb, tendance_data, ...)
    score_churn = models.FloatField(default=0)
    churn_predit = models.BooleanField(default=False)
    niveau_risque = models.CharField(
        max_length=10,
        choices=[("faible", "Faible"), ("eleve", "Élevé")],
        default="faible",
    )

    def save(self, *args, **kwargs):
        if not self.agence_id and self.dataset_id:
            self.agence_id = self.dataset.agence_id
        self.churn_predit = self.calculer_churn()
        super().save(*args, **kwargs)
\end{lstlisting}

Le modèle \texttt{Dataset} encapsule la provenance des données (\texttt{csv}, \texttt{mock}, etc.), le fichier éventuel, l'agence et l'utilisateur ayant chargé le fichier, ainsi que \texttt{nb\_clients}. Lors de l'import CSV, nous créons ou remplaçons des clients liés à ce dataset : la conception distingue ainsi clairement « jeu de données » et « lignes client », ce qui facilite l'historique des analyses par agence.

\section{Flux d'analyse et intégration avec le reste du système}

Lorsqu'un utilisateur lance une analyse avec un fichier, nous lisons le tabulaire, nous peuplons \texttt{Dataset} et \texttt{ClientChurn}, puis nous tentons d'abord un appel à l'API FastAPI (santé via \texttt{/health}, prédiction batch via \texttt{/api/predict/batch}) en reconstruisant les features attendues à partir des champs du modèle Django (par exemple conversion Mo vers Go pour les besoins du pipeline ML). Si l'API échoue ou renvoie un format inattendu, nous rebasculons sur le service ML local. Cette stratégie de repli est une décision de conception explicite : elle nous assure une démonstration stable en environnement étudiant tout en préparant un découplage production (API seule).

\section{Synthèse}

En résumé, nous avons conçu un noyau métier où l'agence structure l'accès aux données, où \texttt{ClientChurn} concentre profil et scores, où les mocks sont câblés dans \path{core/mock_data.py} avec des garde-fous de sécurité, et où l'import CSV et l'API FastAPI s'enchaînent sans confondre données simulées et données réelles. Les diagrammes (diagramme de séquence pour « lancer analyse », schéma entité-association Agence—Dataset—ClientChurn) peuvent être ajoutés dans le rapport maître sous forme de figures ; nous ne les redéfinissons pas ici afin de respecter la mise en page globale du mémoire.

% Si votre classe est article : remplacer \chapter par \section{Conception}
% Nécessite : \usepackage{listings} (ou minted)

\chapter{Conception}

Dans ce chapitre, nous exposons la conception de la partie applicative et péristaltique du système de détection du churn : la manière dont nous structurons les données, la génération contrôlée de \emph{mock} pour le développement et les tests, le modèle métier \texttt{ClientChurn}, et l'articulation avec le reste de la chaîne (datasets réels, API de prédiction, interface). Nous adoptons un point de vue volontairement ancré dans le code : le professeur a insisté sur la présence d'extraits représentatifs ; nous en présentons donc plusieurs, tirés du dépôt tel qu'il existe aujourd'hui.

\section{Périmètre et objectifs de conception}

Nous concevons le système pour qu'un utilisateur rattaché à une \textbf{agence} ne voie et ne modifie que les clients et analyses de cette agence, conformément à la hiérarchie métier. Nous séparons l'interface web (Django), le calcul de scores côté service dédié (FastAPI) lorsqu'il est disponible, et la persistance relationnelle (PostgreSQL via Django). Nous prévoyons en outre un repli sur un scoring local lorsque l'API est injoignable, afin de ne pas bloquer une démonstration ou un environnement de test.

\section{Données mock : rôle, garde-fous et extrait de code}

Les données \emph{mock} nous permettent de peupler rapidement la base avec des profils clients cohérents (télécom, réclamations, usage) lorsque aucun fichier CSV réel n'est chargé, ou pour valider le pipeline bout en bout. Nous avons toutefois restreint explicitement ce mode : si l'application tourne avec \texttt{DEBUG=False} (production), nous interdisons l'appel à la génération mock pour éviter d'injecter des données fictives en situation réelle. Le principe est implémenté dans \path{core/mock_data.py} : une vérification au moment de l'exécution, et une levée d'exception si le mode n'est pas autorisé.

\begin{lstlisting}[language=Python, caption={Garde-fou avant génération mock (\path{core/mock_data.py}, extrait).}]
def _ensure_mock_mode():
    """Leve une exception si le mode Mock est utilise en production."""
    if not _check_mock_mode():
        raise RuntimeError(
            "ERREUR CRITIQUE: Le mode Mock (generer_mock_data) est desactive "
            "en production. DEBUG=False detecte."
        )
\end{lstlisting}

L'entrée principale \texttt{generer\_mock\_data(agence\_id=..., user\_id=..., nb\_clients=50)} résout l'agence et l'utilisateur, crée un \texttt{Dataset} avec \texttt{methode="mock"}, puis génère des enregistrements \texttt{ClientChurn} avec des tirages aléatoires contrôlés (plans tarifaires, CDR agrégés, règles d'audit RGS-90 comme la division protégée pour la durée moyenne d'appel). Nous illustrons ci-dessous la signature et le début de la fonction : c'est le cœur du branchement « données simulées » attendu en conception.

\begin{lstlisting}[language=Python, caption={Entrée de la génération mock (extrait).}]
def generer_mock_data(agence_id=None, user_id=None, nb_clients=50):
    """
    Genere des donnees mock pour les tests.
    SECURITE: bloquee en production (DEBUG=False).
    """
    _ensure_mock_mode()
    from learning.models import ClientChurn, Dataset
    from core.models import Agence
    # ... resolution agence / utilisateur, Dataset(methode="mock"), boucle clients
\end{lstlisting}

Par ailleurs, nous structurons des règles métier de recommandation sous forme de liste de dictionnaires (\texttt{REGLES\_RECOMMANDATIONS}) : chaque entrée combine une condition sur un objet \texttt{ClientChurn} (par exemple nombre de réclamations) et un texte de mission généré dynamiquement. Cette conception nous permet de documenter clairement le lien entre attributs du client mock ou réel et actions proposées dans les notifications, sans durcir ces textes dans la vue.

\section{Modèle \texttt{ClientChurn} et lien \texttt{Dataset} / agence}

Nous modélisons le client comme une entité unique \texttt{ClientChurn} portant à la fois les informations de profil (identité, contrat, usage agrégé) et les sorties du modèle prédictif (\texttt{score\_churn}, \texttt{churn\_predit}, \texttt{niveau\_risque}). Conformément aux migrations du projet, chaque client est rattaché à une \textbf{agence} par clé étrangère, et optionnellement à un \texttt{Dataset} lorsque les lignes proviennent d'un import (CSV) ou d'un jeu mock. Nous complétons l'agence à la sauvegarde si elle manque mais qu'un dataset est présent, pour respecter la contrainte de base et les filtres par agence dans les vues.

\begin{lstlisting}[language=Python, caption={Structure essentielle de \texttt{ClientChurn} (\path{learning/models.py}, extrait).}]
class ClientChurn(models.Model):
    agence = models.ForeignKey("core.Agence", on_delete=models.CASCADE)
    dataset = models.ForeignKey(
        "Dataset", on_delete=models.CASCADE,
        related_name="clients", null=True, blank=True,
    )
    client_id = models.CharField(max_length=100)
    # ... profil, usage (nb_appels, data_totale_mb, tendance_data, ...)
    score_churn = models.FloatField(default=0)
    churn_predit = models.BooleanField(default=False)
    niveau_risque = models.CharField(
        max_length=10,
        choices=[("faible", "Faible"), ("eleve", "Élevé")],
        default="faible",
    )

    def save(self, *args, **kwargs):
        if not self.agence_id and self.dataset_id:
            self.agence_id = self.dataset.agence_id
        self.churn_predit = self.calculer_churn()
        super().save(*args, **kwargs)
\end{lstlisting}

Le modèle \texttt{Dataset} encapsule la provenance des données (\texttt{csv}, \texttt{mock}, etc.), le fichier éventuel, l'agence et l'utilisateur ayant chargé le fichier, ainsi que \texttt{nb\_clients}. Lors de l'import CSV, nous créons ou remplaçons des clients liés à ce dataset : la conception distingue ainsi clairement « jeu de données » et « lignes client », ce qui facilite l'historique des analyses par agence.

\section{Flux d'analyse et intégration avec le reste du système}

Lorsqu'un utilisateur lance une analyse avec un fichier, nous lisons le tabulaire, nous peuplons \texttt{Dataset} et \texttt{ClientChurn}, puis nous tentons d'abord un appel à l'API FastAPI (santé via \texttt{/health}, prédiction batch via \texttt{/api/predict/batch}) en reconstruisant les features attendues à partir des champs du modèle Django (par exemple conversion Mo vers Go pour les besoins du pipeline ML). Si l'API échoue ou renvoie un format inattendu, nous rebasculons sur le service ML local. Cette stratégie de repli est une décision de conception explicite : elle nous assure une démonstration stable en environnement étudiant tout en préparant un découplage production (API seule).

\section{Synthèse}

En résumé, nous avons conçu un noyau métier où l'agence structure l'accès aux données, où \texttt{ClientChurn} concentre profil et scores, où les mocks sont câblés dans \path{core/mock_data.py} avec des garde-fous de sécurité, et où l'import CSV et l'API FastAPI s'enchaînent sans confondre données simulées et données réelles. Les diagrammes (diagramme de séquence pour « lancer analyse », schéma entité-association Agence—Dataset—ClientChurn) peuvent être ajoutés dans le rapport maître sous forme de figures ; nous ne les redéfinissons pas ici afin de respecter la mise en page globale du mémoire.

% Si votre classe est article : remplacer \chapter par \section{Conception}
% Nécessite : \usepackage{listings} (ou minted)

\chapter{Conception}

Dans ce chapitre, nous exposons la conception de la partie applicative et péristaltique du système de détection du churn : la manière dont nous structurons les données, la génération contrôlée de \emph{mock} pour le développement et les tests, le modèle métier \texttt{ClientChurn}, et l'articulation avec le reste de la chaîne (datasets réels, API de prédiction, interface). Nous adoptons un point de vue volontairement ancré dans le code : le professeur a insisté sur la présence d'extraits représentatifs ; nous en présentons donc plusieurs, tirés du dépôt tel qu'il existe aujourd'hui.

\section{Périmètre et objectifs de conception}

Nous concevons le système pour qu'un utilisateur rattaché à une \textbf{agence} ne voie et ne modifie que les clients et analyses de cette agence, conformément à la hiérarchie métier. Nous séparons l'interface web (Django), le calcul de scores côté service dédié (FastAPI) lorsqu'il est disponible, et la persistance relationnelle (PostgreSQL via Django). Nous prévoyons en outre un repli sur un scoring local lorsque l'API est injoignable, afin de ne pas bloquer une démonstration ou un environnement de test.

\section{Données mock : rôle, garde-fous et extrait de code}

Les données \emph{mock} nous permettent de peupler rapidement la base avec des profils clients cohérents (télécom, réclamations, usage) lorsque aucun fichier CSV réel n'est chargé, ou pour valider le pipeline bout en bout. Nous avons toutefois restreint explicitement ce mode : si l'application tourne avec \texttt{DEBUG=False} (production), nous interdisons l'appel à la génération mock pour éviter d'injecter des données fictives en situation réelle. Le principe est implémenté dans \path{core/mock_data.py} : une vérification au moment de l'exécution, et une levée d'exception si le mode n'est pas autorisé.

\begin{lstlisting}[language=Python, caption={Garde-fou avant génération mock (\path{core/mock_data.py}, extrait).}]
def _ensure_mock_mode():
    """Leve une exception si le mode Mock est utilise en production."""
    if not _check_mock_mode():
        raise RuntimeError(
            "ERREUR CRITIQUE: Le mode Mock (generer_mock_data) est desactive "
            "en production. DEBUG=False detecte."
        )
\end{lstlisting}

L'entrée principale \texttt{generer\_mock\_data(agence\_id=..., user\_id=..., nb\_clients=50)} résout l'agence et l'utilisateur, crée un \texttt{Dataset} avec \texttt{methode="mock"}, puis génère des enregistrements \texttt{ClientChurn} avec des tirages aléatoires contrôlés (plans tarifaires, CDR agrégés, règles d'audit RGS-90 comme la division protégée pour la durée moyenne d'appel). Nous illustrons ci-dessous la signature et le début de la fonction : c'est le cœur du branchement « données simulées » attendu en conception.

\begin{lstlisting}[language=Python, caption={Entrée de la génération mock (extrait).}]
def generer_mock_data(agence_id=None, user_id=None, nb_clients=50):
    """
    Genere des donnees mock pour les tests.
    SECURITE: bloquee en production (DEBUG=False).
    """
    _ensure_mock_mode()
    from learning.models import ClientChurn, Dataset
    from core.models import Agence
    # ... resolution agence / utilisateur, Dataset(methode="mock"), boucle clients
\end{lstlisting}

Par ailleurs, nous structurons des règles métier de recommandation sous forme de liste de dictionnaires (\texttt{REGLES\_RECOMMANDATIONS}) : chaque entrée combine une condition sur un objet \texttt{ClientChurn} (par exemple nombre de réclamations) et un texte de mission généré dynamiquement. Cette conception nous permet de documenter clairement le lien entre attributs du client mock ou réel et actions proposées dans les notifications, sans durcir ces textes dans la vue.

\section{Modèle \texttt{ClientChurn} et lien \texttt{Dataset} / agence}

Nous modélisons le client comme une entité unique \texttt{ClientChurn} portant à la fois les informations de profil (identité, contrat, usage agrégé) et les sorties du modèle prédictif (\texttt{score\_churn}, \texttt{churn\_predit}, \texttt{niveau\_risque}). Conformément aux migrations du projet, chaque client est rattaché à une \textbf{agence} par clé étrangère, et optionnellement à un \texttt{Dataset} lorsque les lignes proviennent d'un import (CSV) ou d'un jeu mock. Nous complétons l'agence à la sauvegarde si elle manque mais qu'un dataset est présent, pour respecter la contrainte de base et les filtres par agence dans les vues.

\begin{lstlisting}[language=Python, caption={Structure essentielle de \texttt{ClientChurn} (\path{learning/models.py}, extrait).}]
class ClientChurn(models.Model):
    agence = models.ForeignKey("core.Agence", on_delete=models.CASCADE)
    dataset = models.ForeignKey(
        "Dataset", on_delete=models.CASCADE,
        related_name="clients", null=True, blank=True,
    )
    client_id = models.CharField(max_length=100)
    # ... profil, usage (nb_appels, data_totale_mb, tendance_data, ...)
    score_churn = models.FloatField(default=0)
    churn_predit = models.BooleanField(default=False)
    niveau_risque = models.CharField(
        max_length=10,
        choices=[("faible", "Faible"), ("eleve", "Élevé")],
        default="faible",
    )

    def save(self, *args, **kwargs):
        if not self.agence_id and self.dataset_id:
            self.agence_id = self.dataset.agence_id
        self.churn_predit = self.calculer_churn()
        super().save(*args, **kwargs)
\end{lstlisting}

Le modèle \texttt{Dataset} encapsule la provenance des données (\texttt{csv}, \texttt{mock}, etc.), le fichier éventuel, l'agence et l'utilisateur ayant chargé le fichier, ainsi que \texttt{nb\_clients}. Lors de l'import CSV, nous créons ou remplaçons des clients liés à ce dataset : la conception distingue ainsi clairement « jeu de données » et « lignes client », ce qui facilite l'historique des analyses par agence.

\section{Flux d'analyse et intégration avec le reste du système}

Lorsqu'un utilisateur lance une analyse avec un fichier, nous lisons le tabulaire, nous peuplons \texttt{Dataset} et \texttt{ClientChurn}, puis nous tentons d'abord un appel à l'API FastAPI (santé via \texttt{/health}, prédiction batch via \texttt{/api/predict/batch}) en reconstruisant les features attendues à partir des champs du modèle Django (par exemple conversion Mo vers Go pour les besoins du pipeline ML). Si l'API échoue ou renvoie un format inattendu, nous rebasculons sur le service ML local. Cette stratégie de repli est une décision de conception explicite : elle nous assure une démonstration stable en environnement étudiant tout en préparant un découplage production (API seule).

\section{Synthèse}

En résumé, nous avons conçu un noyau métier où l'agence structure l'accès aux données, où \texttt{ClientChurn} concentre profil et scores, où les mocks sont câblés dans \path{core/mock_data.py} avec des garde-fous de sécurité, et où l'import CSV et l'API FastAPI s'enchaînent sans confondre données simulées et données réelles. Les diagrammes (diagramme de séquence pour « lancer analyse », schéma entité-association Agence—Dataset—ClientChurn) peuvent être ajoutés dans le rapport maître sous forme de figures ; nous ne les redéfinissons pas ici afin de respecter la mise en page globale du mémoire.

% Si votre classe est article : remplacer \chapter par \section{Conception}
% Nécessite : \usepackage{listings} (ou minted)

\chapter{Conception}

Dans ce chapitre, nous exposons la conception de la partie applicative et péristaltique du système de détection du churn : la manière dont nous structurons les données, la génération contrôlée de \emph{mock} pour le développement et les tests, le modèle métier \texttt{ClientChurn}, et l'articulation avec le reste de la chaîne (datasets réels, API de prédiction, interface). Nous adoptons un point de vue volontairement ancré dans le code : le professeur a insisté sur la présence d'extraits représentatifs ; nous en présentons donc plusieurs, tirés du dépôt tel qu'il existe aujourd'hui.

\section{Périmètre et objectifs de conception}

Nous concevons le système pour qu'un utilisateur rattaché à une \textbf{agence} ne voie et ne modifie que les clients et analyses de cette agence, conformément à la hiérarchie métier. Nous séparons l'interface web (Django), le calcul de scores côté service dédié (FastAPI) lorsqu'il est disponible, et la persistance relationnelle (PostgreSQL via Django). Nous prévoyons en outre un repli sur un scoring local lorsque l'API est injoignable, afin de ne pas bloquer une démonstration ou un environnement de test.

\section{Données mock : rôle, garde-fous et extrait de code}

Les données \emph{mock} nous permettent de peupler rapidement la base avec des profils clients cohérents (télécom, réclamations, usage) lorsque aucun fichier CSV réel n'est chargé, ou pour valider le pipeline bout en bout. Nous avons toutefois restreint explicitement ce mode : si l'application tourne avec \texttt{DEBUG=False} (production), nous interdisons l'appel à la génération mock pour éviter d'injecter des données fictives en situation réelle. Le principe est implémenté dans \path{core/mock_data.py} : une vérification au moment de l'exécution, et une levée d'exception si le mode n'est pas autorisé.

\begin{lstlisting}[language=Python, caption={Garde-fou avant génération mock (\path{core/mock_data.py}, extrait).}]
def _ensure_mock_mode():
    """Leve une exception si le mode Mock est utilise en production."""
    if not _check_mock_mode():
        raise RuntimeError(
            "ERREUR CRITIQUE: Le mode Mock (generer_mock_data) est desactive "
            "en production. DEBUG=False detecte."
        )
\end{lstlisting}

L'entrée principale \texttt{generer\_mock\_data(agence\_id=..., user\_id=..., nb\_clients=50)} résout l'agence et l'utilisateur, crée un \texttt{Dataset} avec \texttt{methode="mock"}, puis génère des enregistrements \texttt{ClientChurn} avec des tirages aléatoires contrôlés (plans tarifaires, CDR agrégés, règles d'audit RGS-90 comme la division protégée pour la durée moyenne d'appel). Nous illustrons ci-dessous la signature et le début de la fonction : c'est le cœur du branchement « données simulées » attendu en conception.

\begin{lstlisting}[language=Python, caption={Entrée de la génération mock (extrait).}]
def generer_mock_data(agence_id=None, user_id=None, nb_clients=50):
    """
    Genere des donnees mock pour les tests.
    SECURITE: bloquee en production (DEBUG=False).
    """
    _ensure_mock_mode()
    from learning.models import ClientChurn, Dataset
    from core.models import Agence
    # ... resolution agence / utilisateur, Dataset(methode="mock"), boucle clients
\end{lstlisting}

Par ailleurs, nous structurons des règles métier de recommandation sous forme de liste de dictionnaires (\texttt{REGLES\_RECOMMANDATIONS}) : chaque entrée combine une condition sur un objet \texttt{ClientChurn} (par exemple nombre de réclamations) et un texte de mission généré dynamiquement. Cette conception nous permet de documenter clairement le lien entre attributs du client mock ou réel et actions proposées dans les notifications, sans durcir ces textes dans la vue.

\section{Modèle \texttt{ClientChurn} et lien \texttt{Dataset} / agence}

Nous modélisons le client comme une entité unique \texttt{ClientChurn} portant à la fois les informations de profil (identité, contrat, usage agrégé) et les sorties du modèle prédictif (\texttt{score\_churn}, \texttt{churn\_predit}, \texttt{niveau\_risque}). Conformément aux migrations du projet, chaque client est rattaché à une \textbf{agence} par clé étrangère, et optionnellement à un \texttt{Dataset} lorsque les lignes proviennent d'un import (CSV) ou d'un jeu mock. Nous complétons l'agence à la sauvegarde si elle manque mais qu'un dataset est présent, pour respecter la contrainte de base et les filtres par agence dans les vues.

\begin{lstlisting}[language=Python, caption={Structure essentielle de \texttt{ClientChurn} (\path{learning/models.py}, extrait).}]
class ClientChurn(models.Model):
    agence = models.ForeignKey("core.Agence", on_delete=models.CASCADE)
    dataset = models.ForeignKey(
        "Dataset", on_delete=models.CASCADE,
        related_name="clients", null=True, blank=True,
    )
    client_id = models.CharField(max_length=100)
    # ... profil, usage (nb_appels, data_totale_mb, tendance_data, ...)
    score_churn = models.FloatField(default=0)
    churn_predit = models.BooleanField(default=False)
    niveau_risque = models.CharField(
        max_length=10,
        choices=[("faible", "Faible"), ("eleve", "Élevé")],
        default="faible",
    )

    def save(self, *args, **kwargs):
        if not self.agence_id and self.dataset_id:
            self.agence_id = self.dataset.agence_id
        self.churn_predit = self.calculer_churn()
        super().save(*args, **kwargs)
\end{lstlisting}

Le modèle \texttt{Dataset} encapsule la provenance des données (\texttt{csv}, \texttt{mock}, etc.), le fichier éventuel, l'agence et l'utilisateur ayant chargé le fichier, ainsi que \texttt{nb\_clients}. Lors de l'import CSV, nous créons ou remplaçons des clients liés à ce dataset : la conception distingue ainsi clairement « jeu de données » et « lignes client », ce qui facilite l'historique des analyses par agence.

\section{Flux d'analyse et intégration avec le reste du système}

Lorsqu'un utilisateur lance une analyse avec un fichier, nous lisons le tabulaire, nous peuplons \texttt{Dataset} et \texttt{ClientChurn}, puis nous tentons d'abord un appel à l'API FastAPI (santé via \texttt{/health}, prédiction batch via \texttt{/api/predict/batch}) en reconstruisant les features attendues à partir des champs du modèle Django (par exemple conversion Mo vers Go pour les besoins du pipeline ML). Si l'API échoue ou renvoie un format inattendu, nous rebasculons sur le service ML local. Cette stratégie de repli est une décision de conception explicite : elle nous assure une démonstration stable en environnement étudiant tout en préparant un découplage production (API seule).

\section{Synthèse}

En résumé, nous avons conçu un noyau métier où l'agence structure l'accès aux données, où \texttt{ClientChurn} concentre profil et scores, où les mocks sont câblés dans \path{core/mock_data.py} avec des garde-fous de sécurité, et où l'import CSV et l'API FastAPI s'enchaînent sans confondre données simulées et données réelles. Les diagrammes (diagramme de séquence pour « lancer analyse », schéma entité-association Agence—Dataset—ClientChurn) peuvent être ajoutés dans le rapport maître sous forme de figures ; nous ne les redéfinissons pas ici afin de respecter la mise en page globale du mémoire.
